%% ma: L10-B14-0001
%% lop: 10
%% chuong: 5
%% bai: 14
%% dang: 10.14.2
%% loai: TL
%% muc: TH
%% dap_an: "-"
%% nguon: "chuyen-de-cac-so-dac-trung-cua-mau-so-lieu-khong-ghep-nhom-toan-10-chuong-trinh-moi.pdf (GV. Phan Nhật Linh) (P04), Bài 03 Dạng 1 Bài tập tự luận 1"
%% kien_thuc: "Lớp 10 Bài 13 (tứ phân vị), Bài 14 (khoảng tứ phân vị)"
%% trang_thai: chinh-thuc
%% ghi_chu: "Tự duyệt theo yêu cầu giáo viên 22/10/2026."
\begin{ex}
Chiều cao (đơn vị: mét) của $15$ cây bạch đàn được thống kê như sau:
\begin{center}
\begin{tabular}{|c|c|c|c|c|c|c|c|}\hline
$6{,}3$ & $6{,}6$ & $7{,}5$ & $8{,}2$ & $8{,}3$ & $7{,}8$ & $7{,}9$ & $9{,}0$\\\hline
$8{,}9$ & $7{,}2$ & $7{,}5$ & $8{,}7$ & $7{,}7$ & $8{,}8$ & $7{,}6$ & \\\hline
\end{tabular}
\end{center}
a) Tìm khoảng biến thiên của mẫu số liệu.

b) Tìm khoảng tứ phân vị của mẫu số liệu.
\loigiai{a) $R=9{,}0-6{,}3=2{,}7$ (m).

b) Sắp xếp: $6{,}3$; $6{,}6$; $7{,}2$; $7{,}5$; $7{,}5$; $7{,}6$; $7{,}7$; $7{,}8$; $7{,}9$; $8{,}2$; $8{,}3$; $8{,}7$; $8{,}8$; $8{,}9$; $9{,}0$.

$Q_2=x_8=7{,}8$; nửa dưới gồm $7$ giá trị đầu nên $Q_1=x_4=7{,}5$; nửa trên gồm $7$ giá trị cuối nên $Q_3=x_{12}=8{,}7$.

$\Delta_Q=8{,}7-7{,}5=1{,}2$ (m).}
\end{ex}
